% Appendix C. Fragment, \input by spacenav.tex inside \appendix.
% Transcribed from spacenav/env.py: _sample_waypoint, sample_tour, current_leg,
% waypoint_states, deltav_estimate, _substep, step; spacenav/physics.py
% hazard_terms; spacenav/constants.py.

\section{Goal conditions}
\label{app:goals}

A task instance is a tour of between one and six waypoints, sampled one after
another by \code{sample\_tour()} in \code{spacenav/env.py} so that each lies at
least 100\,u (\code{min\_leg}) from the one before it, and each is of one of
the three kinds in Table~\ref{tab:goal-kinds}, which differ in what the
waypoint is attached to and therefore in how it moves. Whether a waypoint
moves is decided first (with probability \code{p\_station}, 0.35 by default
and set per suite), and a moving waypoint is a station when the level has a
free one and a coin comes up heads, otherwise a point hung off a solid body; a
waypoint that does not move is a point sampled uniformly over the playable
disc, rejected and redrawn (up to 96 candidates) until it lies between
\code{min\_dist} and \code{max\_dist} from the previous waypoint, inside 0.9 of
the level radius, at least 20\,u clear of every surface and atmosphere, and at
a point where the instantaneous damage rate is under 0.3 health/s. The
velocity condition is drawn second and only for a moving waypoint, with
probability \code{p\_rendezvous} (so that the rendezvous shares quoted in
Table~\ref{tab:suites} are conditional on the waypoint moving, and a fixed
point in free space is always a flythrough). Every waypoint has the same
arrival radius of 12\,u (\code{TARGET\_RADIUS}), and the craft is at a waypoint
when, at the end of a physics substep, both its distance to the waypoint's
current position is under 12\,u and its speed relative to the waypoint's
current velocity is under the waypoint's tolerance, which is 3\,u/s
(\code{RENDEZVOUS\_V\_TOL}) for a rendezvous and $10^9$\,u/s, i.e. any speed,
for a flythrough. Both conditions are evaluated against the waypoint's state
at that substep, so a rendezvous with a station is a velocity match with the
station's orbital motion and a rendezvous with a point hung off a giant is a
match with the giant's.

\begin{table}[htbp]
\centering
\footnotesize
\begin{tabular}{p{2.5cm}p{4.6cm}p{2.7cm}p{2.9cm}}
\toprule
Kind & Where it is & How it moves & Velocity condition \\
\midrule
\raggedright free-space point & \raggedright a fixed position in the level frame, sampled as above (\code{wp\_anchor} $= -1$) & \raggedright it does not; its velocity is zero & \raggedright flythrough only \tabularnewline
\raggedright station & \raggedright the position of a station body (massless, zero radius, on its own orbit around its host) & \raggedright with the station & \raggedright flythrough or rendezvous with the station's velocity \tabularnewline
\raggedright attached to a body & \raggedright the body's position plus a constant offset vector of length (radius + atmosphere + 35\,u) at a random angle, or, for a contested waypoint, of length $0.8$--$1.3\times$ the body's equal-pull radius (Section~\ref{sec:eta}) & \raggedright the offset translates with the body and does not rotate with it & \raggedright flythrough or rendezvous with the body's velocity \tabularnewline
\bottomrule
\end{tabular}
\caption{The three waypoint kinds. Every kind has the same 12\,u arrival
radius; black holes and tracers are never anchors, and a body that already
carries a waypoint is not chosen again in the same tour. The figure of the
task (Figure~\ref{fig:task}d) shows the same conditions as four cells by
separating the velocity condition from the attachment.}
\label{tab:goal-kinds}
\end{table}

A tour is flown in slot order unless the instance's \code{order\_free} flag is
set (with probability \code{p\_order\_free}, zero in training until the
curriculum raises it and 0.3 on the blind suite), and the difference between
the two is only in which waypoint an arrival is credited to: in a fixed-order
tour, reaching a waypoint counts only if it is the next unvisited one in slot
order, so that passing through the third waypoint on the way to the first does
nothing, while in an order-free tour any unvisited waypoint counts whenever its
conditions are met. The observation's notion of the current target
(Appendix~\ref{app:obs}) follows the same rule, being the next unvisited
waypoint in a fixed-order tour and the nearest unvisited one in an order-free
tour, and so does the potential used for reward shaping, which is the negative
of the distance to the current target plus a flat allowance for each further
waypoint still outstanding. The scoring of a completed tour is identical in
the two cases (the same per-waypoint and completion rewards, the same weighted
cost over time, path, fuel and exposure), so an order-free tour is neither
easier nor harder by construction; it is different in that the navigator, and
not the instance, chooses the sequence.

The craft's loadout is drawn against the tour it has to fly. With $d_k$ the
straight-line length of the $k$-th leg (from the start position to the first
waypoint and then from each waypoint to the next, in slot order and ignoring
gravity, the craft's initial velocity and any motion of the waypoints), the
$\Delta v$ the tour is estimated to require is
\begin{equation}
\Delta v_{\mathrm{req}} \;=\; \sum_k
\begin{cases}
\sqrt{2\,a_{\mathrm{eng}}\,d_k} & \text{flythrough, i.e. accelerate across the whole leg,}\\[2pt]
2\sqrt{a_{\mathrm{eng}}\,d_k} & \text{rendezvous, i.e. accelerate for half the leg and brake for half,}
\end{cases}
\end{equation}
which \code{deltav\_estimate()} computes and which is also what the
observation reports as the $\Delta v$ required over the legs still pending.
The engine acceleration $a_{\mathrm{eng}}$ is drawn log-uniformly from
$[0.6, 3]$\,u/s$^2$ (\code{SHIP\_ACCEL\_RANGE}, narrowed to $[1.5, 3]$ on the
easy suite and $[0.8, 2]$ on the hard suite, and started generous by the
training curriculum), the budget multiple $\beta$ is drawn uniformly from
$[0.55, 1.35]$ (\code{DELTAV\_BUDGET\_RANGE}, narrowed per suite as
Table~\ref{tab:suites} lists), and the tank, in seconds of full-throttle burn,
is $\beta\,\Delta v_{\mathrm{req}} / a_{\mathrm{eng}}$ clipped to
$[3, 60]$\,s (\code{SHIP\_FUEL\_LIMITS}), so that the tank actually carried can
be a larger multiple of the estimate than $\beta$ on a very short tour and a
smaller one on a very long tour. The reference scales that normalize the four
cost terms are set from the same quantities, with the time and exposure
references equal to $\sum_k 2\sqrt{d_k/a_{\mathrm{eng}}}$, the path reference
equal to $\sum_k d_k$ and the fuel reference equal to half the tank, and the
objective weights are either a single objective (with probability 0.5, drawn
with probabilities 0.4, 0.2, 0.3 and 0.1 over time, path, fuel and exposure)
or a Dirichlet mixture of them, with exposure available only on a level that
has at least one zone, the zone charged for being drawn from the level's
active zones with sensor zones favored by a factor of $e$ over the rest.

An episode ends in exactly one of five ways, tested in this order at the end
of every physics substep (60\,Hz) and reported as the status code the
evaluation suites count: \code{crashed}, the craft's position is inside the
radius of a solid body (a station or a tracer, having no radius, cannot be
crashed into); \code{destroyed}, health has reached zero under the accumulated
damage from stellar radiation, atmospheric heating, debris and radiation
belts; \code{arrived}, every active waypoint has been visited, which is the
only success; \code{lost}, the craft is farther than 1.35 times the level
radius from the origin; and, tested once per control tick rather than per
substep, \code{timeout}, the tick count has reached $T_{\max}/\Delta t$ (2250
at the 150\,s clock) with the tour unfinished. An empty tank is not among
them. The burn in a substep is the smaller of the commanded burn and the fuel
left, so that once the tank is empty the throttle does nothing, the craft
coasts under gravity alone, and the episode continues until one of the five
endings occurs, which on an energy-efficient path is often the clock
(Section~\ref{sec:clock}) and on a decaying one is a crash. The hazard terms
themselves do not end an episode; they lower health, and only health reaching
zero does.
